\documentclass[11pt]{article}
\usepackage[margin=1.02in]{geometry}
\usepackage[T1]{fontenc}
\usepackage{lmodern,amsmath,amssymb,amsthm,microtype}
\usepackage[hidelinks]{hyperref}
\hypersetup{pdftitle={The abundance of weird numbers of the form 2\string^k pq},pdfauthor={Leo Gao}}
\usepackage{enumitem}
\setlist[itemize]{leftmargin=1.25em,itemsep=0.45em,topsep=0.4em}
\newtheorem*{theorem}{Theorem}
\newtheorem*{corollary}{Corollary}
\setlength{\parskip}{0.25em}
\title{The abundance of weird numbers of the form \(2^kpq\)}
\author{Leo Gao}
\date{27 November 2026}
\begin{document}
\maketitle
\begin{abstract}
We prove Iannucci's conjectured upper bound for the abundance of a weird number of the form \(2^kpq\). If the bound fails, a modular progression must remain in an interval shorter than half its modulus. Its common difference must then be one of eight integers; none permits weirdness. As a consequence, \(70\) uniquely maximizes the abundancy index in this family.
\end{abstract}

An integer \(n\) is \emph{weird} if its proper divisors sum to more than \(n\), but no subset sums to \(n\). Write \(\sigma(n)=\sum_{d\mid n}d\). Benkoski and Erd\H os \cite{BE} asked whether there are infinitely many primitive weird numbers, those with no proper weird divisor. The family \(2^kpq\), with \(p<q\) odd primes, has played a central role: its weird members are primitive, and Melfi \cite{M} gives a conditional infinitude result. Here we prove a different assertion about this family, conjectured from numerical evidence by Iannucci \cite[\S III, (6)]{I}. It does not settle infinitude.

\begin{theorem}
Let \(n=2^kpq\) be weird, where \(k\ge1\) and \(p<q\) are odd primes. For \(M=2^{k+1}-1\) and \(D=\sigma(n)-2n\),
\[
                 M<D<\frac{(M+1)(M+2)}{3}.
\]
\end{theorem}
\begin{proof}
The sum of all proper divisors is \(n+D\). By complementation, weirdness implies that \(D>0\) is not a sum of distinct proper divisors. Since every integer from 0 to \(M\) is a subset sum of the powers of two \(1,2,4,\ldots,2^k\), \(D>M\). Multiplicativity gives
\[
 D=M+Mp+Mq-pq=M(M+1)-(p-M)(q-M).
\]
Any representation \(pq=r+bp+cq\) with integers \(0\le r,b,c\le M\) contradicts weirdness: represent the three coefficients in binary using the disjoint divisor blocks \(2^i,2^ip,2^iq\) with \(0\le i\le k\), and add \((2^k-1)pq\) using all \(2^ipq\) with \(0\le i<k\). In particular \(p>M\), since otherwise \(pq=p\cdot q\) uses coefficient \(c=p\). From \(D>M\) we get \((p-M)(q-M)<M^2\), hence \(p<2M\). Furthermore \(q\ge 2M+1\): if \(q\le 2M-1\), the coefficients in
\[
 pq=\tfrac{q-p}{2}+\tfrac{q+1}{2}p+\tfrac{p-1}{2}q
\]
belong to \([0,M]\).

Put \(d=p-M\), a positive even integer. Suppose to the contrary
\[
 d(q-M)\le T:=2(M^2-1)/3. \tag{1}
\]
Since \(q-M\ge M+1\), this implies \(d\le 2(M-1)/3<2M/3\). For \(M=3\) there is no such positive even \(d\); hence \(M\ge7\).

Let \(A=\lfloor M(p+1)/q\rfloor\) and \(h=A-d\). For any \(a\in[d,A]\), divide \(aq=bp+r\) with \(0\le r<p\). We have \(b\le M\) because \(aq\le M(p+1)<(M+1)p\). Also \(a\le p\) since \(q>2M\), so \(0\le p-a\le M\). If \(r\le M\), then \(pq=r+bp+(p-a)q\) is a forbidden certificate. Therefore for every such \(a\) the residues
\[
 z_a=(-aq\bmod p)\quad\text{belong to }[1,d-1]. \tag{2}
\]
\paragraph{A short modular progression.} We have
\[
 \frac{M(p+1)}q-d\ \ge \frac{d(M^2+3M+2)}{2M^2+3Md-2}>d/4,
\]
using (1) and \(3d<2M\). For even \(d\ge4\), this gives \(h\ge\lfloor d/4\rfloor\ge(d-2)/4\) and \(h\ge1\). If \(d=2\), (1) gives \(3q\le M^2+3M-1<M(p+1)\), hence \(h\ge1\); (2) makes two successive residues identical, impossible since \(p\nmid q\). Now take \(d\ge4\). Successive differences \(z_{a+1}-z_a\) lie in \([-(d-2),d-2]\), an interval of length less than \(p=M+d\), and all are congruent to \(-q\bmod p\). They are therefore the same nonzero integer \(w\), with \(h|w|\le d-2\) and so \(1\le |w|\le4\). In particular,
\[
 q=jp-w\quad(j\in\mathbb Z),\qquad z_d=dw\bmod p\in[1,d-1]. \tag{3}
\]

\paragraph{Eliminating the slopes.} Throughout \(d<M\), \(p=M+d\), and \(d<2M/3\).

\begin{itemize}
\item \(w=1,-1,2\): \(z_d\) is respectively \(d,M,2d\), none admissible.
\item \(w=-2\): \(z_d=M-d<d\) requires \(d>M/2\). As \(p+2<2M+1\) (since \(M\ge7\) and \(d<2M/3\)), parity in (3) gives \(q\ge3p+2\). Then \(d(q-M)>M^2>T\).
\item \(w=3\): \(3d<2p\); the only admissible residue is \(3d-p\), requiring \(d>M/2\), hence \(d\ge(M+1)/2\). Parity gives \(q\ge2p-3\). Thus \(d(q-M)\ge (M+1)(2M-2)/2=M^2-1>T\).
\item \(w=-3\): if \(3d\ge p\), the residue is either zero or \(2p-3d=2M-d>d\). Otherwise (3) gives \(z_d=M-2d\in(0,d)\), so \(M/3<d<M/2\). Parity gives either \(q\ge4p+3\), contradicting (1) since \(d>M/3\), or \(q=2p+3\). In the latter case \(2d\le M-3\): the only other even integer \(2d<M\) is \(M-1\), which gives \(d(q-M)= (M-1)(M+1)>T\). Set \(a=\lceil p/3\rceil\). We have \(a\ge d\) and \(a\le (M-1)/2\) since \(2d\le M-3\). Thus
\[
 pq=(3a-p)+(2a+1)p+(p-a)q
\]
has three coefficients in \([0,M]\), contradicting weirdness.
\item \(w=4\): since \(4d<2p\), (3) requires \(4d>p\), i.e. \(d>M/3\). Parity and \(q>p\) give \(q\ge3p-4\). Thus \(d(q-M)>(M/3)(3M-4)>T\) for \(M\ge7\).
\item \(w=-4\): if \(4d\ge p\), the residue is either zero or \(2p-4d=2M-2d>d\). Otherwise (3) gives \(M/4<d<M/3\). Now \(q=p+4\) is impossible since \(p+4<2M+1\) for \(M\ge7\) and \(d<M/3\); parity forces \(q\ge3p+4\). Therefore \(d(q-M)>(M/4)(2M+3M/4)=11M^2/16>T\).

\end{itemize}
All possibilities contradict (1) or weirdness. Thus \(d(q-M)>T\), which is precisely the desired strict upper bound.
\end{proof}

\clearpage
\begin{corollary}
Among weird numbers of the form \(2^kpq\), the largest abundancy index \(\sigma(n)/n\) is \(72/35\), attained only at \(70\).
\end{corollary}
\begin{proof}
The theorem and its proof give \(p\ge M+2\) and \(q\ge2M+1\). Hence
\[
\frac{D}{n}<\frac{2(M+2)}{3pq}\le\frac{2}{3(2M+1)}.
\]
For \(k\ge2\) the last expression is at most \(2/45<2/35\). For \(k=1\), the conditions \(3<p<6\) and \((p-3)(q-3)<9\) force \((p,q)=(5,7)\). The abundance of 70 is 4, not attainable from its proper divisors: only 1 and 2 are smaller than 5. Thus 70 is weird and its abundancy index is \(2+2/35\).
\end{proof}

\paragraph{Research note.} This manuscript was prepared with AI assistance. No part of the proof depends on computation. The accompanying optional script independently tests the statement by direct subset sums on finite data.

\begin{thebibliography}{9}
\bibitem{BE} S. J. Benkoski and P. Erd\H os, \href{https://doi.org/10.1090/S0025-5718-1974-0347726-9}{On weird and pseudoperfect numbers}, \emph{Math. Comp.} \textbf{28} (1974), 617--623.
\bibitem{I} D. E. Iannucci, \href{https://arxiv.org/abs/1504.02761}{On primitive weird numbers of the form \(2^kpq\)}, \emph{IOSR J. Math.} \textbf{11}, no. 3, ver. III (2015), 86--130.
\bibitem{M} G. Melfi, \href{https://doi.org/10.1016/j.jnt.2014.07.024}{On the conditional infiniteness of primitive weird numbers}, \emph{J. Number Theory} \textbf{147} (2015), 508--514.
\end{thebibliography}
\end{document}
